\documentclass[11pt,a4paper]{article}
\usepackage[margin=2.2cm]{geometry}
\usepackage[utf8]{inputenc}\usepackage[T1]{fontenc}
\usepackage{amsmath,amssymb}
\usepackage{fancyhdr}\usepackage{enumitem}\usepackage{xcolor}
\usepackage{hyperref}
\pagestyle{fancy}\fancyhf{}
\rhead{Math 1A Worksheet 3}
\lhead{Songpengfei}
\rfoot{Page \thepage}
\newcommand{\solvertag}[1]{{\footnotesize\color{gray}\textsf{[#1]}}}

\title{Homework Solutions}
\author{Songpengfei}
\date{2026-10-09}
\begin{document}
\maketitle

\noindent\textbf{Solved: 8/8 (100.0\%)}\quad\textit{Auto-generated by C4C homework-solver.}

\section*{Problem Q1}
Can the graph of a function have more than one tangent at a given point? If so, graph your answer. If not, explain why.

\paragraph{Solution.} \solvertag{conceptual\_template} 

\begin{enumerate}[leftmargin=*]
  \item In general, a smooth function has exactly one tangent line at each point.
  \item However, a curve (not necessarily the graph of a function) can have multiple tangents at a self-intersection point.
  \item For a function graph, the tangent is unique where the derivative exists.
\end{enumerate}

\noindent\textbf{Answer:}\quad $\text{Unique tangent where } f'(a) \text{ exists.}$

\medskip

\section*{Problem Q2}
Is there a function whose graph doesn't have a tangent at some point? If so, graph your answer. If not, explain why.

\paragraph{Solution.} \solvertag{conceptual\_template} 

\begin{enumerate}[leftmargin=*]
  \item Yes. A function may fail to have a tangent at a point where it is not differentiable.
  \item Examples: $f(x) = |x|$ at $x = 0$ (corner), $f(x) = x^{1/3}$ at $x = 0$ (vertical tangent/cusp).
\end{enumerate}

\noindent\textbf{Answer:}\quad $\text{Yes: } f(x) = |x| \text{ at } x = 0 \text{ (corner)}$

\medskip

\section*{Problem P1}
Suppose $\delta$ is a positive real number. How do you describe all real numbers $x$ that are within $\delta$ of 0?

\paragraph{Solution.} \solvertag{epsilon\_delta\_template} 

\begin{enumerate}[leftmargin=*]
  \item x within δ of 0: $|x| < \delta$, i.e. $-\delta < x < \delta$
\end{enumerate}

\paragraph{(a)} $-\delta < x < \delta$

\paragraph{(b)} $|x| < \delta$

\paragraph{(c)} $(-\delta, \delta)$

\noindent\textbf{Answer:}\quad $|x| < \delta$, i.e. $-\delta < x < \delta$

\medskip

\section*{Problem P2}
Suppose $a$ is a real number and $\delta$ is a positive real number. How do you describe all real numbers $x$ that are within $\delta$ of $a$?

\paragraph{Solution.} \solvertag{epsilon\_delta\_template} 

\begin{enumerate}[leftmargin=*]
  \item x within δ of a: $|x - a| < \delta$, i.e. $a - \delta < x < a + \delta$
\end{enumerate}

\paragraph{(a)} $\text{See graph}$

\paragraph{(b)} $a - \delta < x < a + \delta$

\paragraph{(c)} $|x - a| < \delta$

\paragraph{(d)} $(a - \delta, a + \delta)$

\noindent\textbf{Answer:}\quad $|x - a| < \delta$, i.e. $a - \delta < x < a + \delta$

\medskip

\section*{Problem P3}
Suppose $a$ is a real number and $\delta$ is a positive real number. How do you describe all real numbers $x$ that are within $\delta$ of $a$, but not equal to $a$:

\paragraph{Solution.} \solvertag{epsilon\_delta\_template} 

\begin{enumerate}[leftmargin=*]
  \item x within δ of a, not equal to a: $0 < |x - a| < \delta$
\end{enumerate}

\paragraph{(a)} $\text{See graph}$

\paragraph{(b)} $a - \delta < x < a + \delta, \quad x \neq a$

\paragraph{(c)} $0 < |x - a| < \delta$

\paragraph{(d)} $(a - \delta, a) \cup (a, a + \delta)$

\noindent\textbf{Answer:}\quad $0 < |x - a| < \delta$

\medskip

\section*{Problem P4}
Let $f(x) = x^2$.

\paragraph{Solution.} \solvertag{epsilon\_delta\_computation} 

\begin{enumerate}[leftmargin=*]
  \item Given: $f(x) = x^{2}$
\end{enumerate}

\paragraph{(a)} $(2.82842712474619, 3.16227766016838)$

\paragraph{(b)} $\delta = 0.16227766016838$

\paragraph{(c)} {\color{red}\textit{unsolved: 无法解析 ε-δ 计算子题}}

\paragraph{(d)} $\text{Yes (continuity)}$

\noindent\textbf{Answer:}\quad $f(x) = x^{2}$

\medskip

\section*{Problem AP1}
Graph $y = 2x - x^2$. For which points $a$ is the tangent line to the curve at the point $(a, 2a - a^2)$ a horizontal line? How is the value of $2x - x^2$ at these points related to the values of $2x - x^2$ at other points?

\paragraph{Solution.} \solvertag{sympy} 

\begin{enumerate}[leftmargin=*]
  \item Given: $y = - x^{2} + 2 x$
  \item Derivative: $y' = 2 - 2 x$
  \item Horizontal tangent where $y' = 0$: $2 - 2 x = 0$
  \item Solutions: $x = \left[ 1\right]$
  \item Points: $(1, 1)$
  \item $f''(1) = -2 < 0$: this is a maximum.
  \item The value $1$ is the maximum of $- x^{2} + 2 x$.
\end{enumerate}

\noindent\textbf{Answer:}\quad $(1, 1)$

\medskip

\section*{Problem AP2}
The point $(2, 0)$ is on the curve $y = 2x - x^2$.

\paragraph{Solution.} \solvertag{sympy} 

\begin{enumerate}[leftmargin=*]
  \item Given: $y = - x^{2} + 2 x$
  \item Derivative: $y' = 2 - 2 x$
  \item At point $(2, 0)$:
  \item Slope: $f'(2) = -2$
  \item Tangent line: $y = 4 - 2 x$
\end{enumerate}

\paragraph{(a)} $\text{Line with slope } -2$

\paragraph{(b)} $y = 4 - 2 x$

\noindent\textbf{Answer:}\quad $4 - 2 x$

\medskip

\end{document}
